%% ch14 --- Weather, climate and the ensemble forecast
%%
%% Stub, generated by tools/gen_stubs.py from tools/chapters.json.
%% The brief below is the contract for this chapter: what it must argue, what
%% it must buy the reader, and what it must NOT take from its neighbours.
%% Writing the chapter means deleting the stub block below (the one that opens
%% on the next \chapter line) and replacing it with the chapter text -- after
%% which this file is never regenerated.
%% If the brief turns out to be wrong once the code has been run, change it in
%% the manifest and record the contradiction in CLAUDE.md under *Resolved
%% questions*.

\chapter{Weather, climate and the ensemble forecast}
\label{ch:weather}

\chapterstub{%
Argument: weather forecasting is the place chaos theory changed practice
most. Because the atmosphere has a positive Lyapunov exponent and many
scales of motion, Lorenz argued in 1969 that detailed weather has a finite
predictability of about two weeks however good the data \dash{} the commonly
quoted figure, labelled as such. Forecasters answered with the ENSEMBLE: run
the model many times from slightly different starts consistent with what was
measured, and report the spread \dash{} a forecast becomes a probability.
The reader builds one on the Lorenz system and on Lorenz's 1996 forty-
variable toy atmosphere: the ensemble spread grows with lead time, and the
spread itself says how far ahead the forecast is worth trusting. Weather
against climate: climate is the statistics of the attractor, which Chapter 9
showed are predictable when individual paths are not, so the claim that we
cannot predict the weather next month so we cannot predict climate confuses
the path with the attractor \dash{} and changing a parameter (the forcing)
moves the statistics, which is what a climate projection asks. Payoff: the
reader can build an ensemble forecast, read its spread as confidence,
compute a probability from it, and separate a weather question from a
climate question. Traps to elicit: a forecast is a single number; if the
weather is unpredictable the climate is too; a 30 per cent chance of rain
forecast that was followed by rain was wrong. Listings: an ensemble on the
Lorenz system; ensemble spread against lead time; Lorenz-96 ensemble;
climate statistics under two parameter values. Labs: turn an ensemble into a
probability; measure the lead time at which spread saturates. Figures: an
ensemble plume fanning out; spread against lead time; statistics shifting
with a parameter.

\medskip
\textbf{Word budget:} 3500.
}
